\appendix
\FloatBarrier
\section{Proofs}\label{app:proofs}

\paragraph{Proof of Proposition~\ref{prop:repeller}.} Throughout, $F_C(S) = \max(S + g(S) - C,\,0)$; the numerical floor $\epsilon = 10^{-3}\kt$ used in the computations replaces $0$ by $\epsilon$ and changes nothing below.

\emph{(i)} Let $0 \le S < S_-(C)$. If $g(S) \ge C$ then $S$ belongs to the set $\{S'>0 : g(S')\ge C\}$ (for $S>0$) or $g(0)=0\ge C>0$ (for $S=0$), both impossible; hence $g(S) < C$ and $S + g(S) - C < S$. Therefore $0 \le F_C(S) \le \max(S,0) = S < S_-(C)$. By induction a trajectory started in $[0,S_-(C))$ is nonincreasing and never leaves the interval. If $g(S) \ge C$ on an interval $I$ and $S\in I$, then $S + g(S) - C \ge S \ge 0$, so $F_C(S) = S + g(S) - C \ge S$.

\emph{(ii)} Suppose $S_{t+1} = F_{C_t}(S_t) < S_t - C_t$. If the maximum in $F_{C_t}$ is attained by the first argument, $S_t + g(S_t) - C_t < S_t - C_t$ gives $g(S_t) < 0$. If it is attained by $0$, then $S_t + g(S_t) - C_t \le 0 < S_t - C_t$ (the right inequality because $S_{t+1} = 0 < S_t - C_t$), and again $g(S_t) < 0$. On the compensatory branch $g(S) = rS(1-S/K) \ge 0$ for $0 \le S \le K$. On the depensatory branch $g(S) = rS(1-S/K)(S-\sA)/(K-\sA)$ with $0<\sA<K$, which for $0<S<K$ is negative exactly when $S<\sA$. \qed

For the Schaefer branch the equation $g(S)=C$ has the roots $S_\pm(C) = \tfrac{K}{2}\bigl(1\pm\sqrt{1-4C/(rK)}\bigr)$ when $4C\le rK$, $g\ge C$ on $[S_-,S_+]$ and $g<C$ outside, so $S_-(C)$ in the proposition is the lower root and $[S_-(C),S_+(C)]$ the interval on which the autonomous map cannot decline. $F_C'(S_+) = 1 + g'(S_+) = 1 + r(1-2S_+/K)$, which at the collapse-window fit under $C=240\kt$ equals $-0.39$.

\paragraph{Proof of Proposition~\ref{prop:increment}.} The forecast error of $\hat y$ is $\Delta_{t+1}-m(\hat x_{t+1}) = (m_{t+1}-m(\hat x_{t+1})) + (\Delta_{t+1}-m_{t+1}) = d_{t+1}+u_{t+1}$, and the error of persistence is $\Delta_{t+1} = m_{t+1}+u_{t+1}$. Hence
\begin{align*}
\MSE(\hat y)-\MSE(\pers) &= \E[d^2]+2\E[du]+\E[u^2] - \E[m^2]-2\E[mu]-\E[u^2]\\
&= \E[d^2]-\E[m^2] + 2\E\bigl[u\,(d-m_{t+1})\bigr],
\end{align*}
and $d-m_{t+1} = -m(\hat x_{t+1})$, which is \eqref{eq:decomp}.

\emph{(i)} With $\hat x_{t+1}=x_{t+1}$, $d\equiv 0$ and \eqref{eq:decomp} reads $-\E[m^2]-2\E[um] = \E[u^2]-\E[(u+m)^2] = \E[u^2]-\E[\Delta^2]$.

\emph{(ii)} With $\hat x_{t+1}=x_t$, $d_{t+1}=m_{t+1}-m_t$, and weak stationarity gives $\E[d^2] = 2\sigma^2 - 2\operatorname{cov}(m_{t+1},m_t) = 2\sigma^2(1-\rho)$ and $\E[m^2]=\sigma^2+\mu^2$. The misspecification term is $-2\E[u_{t+1}m_t]=0$ by hypothesis. Hence $\MSE(\hat y)-\MSE(\pers) = 2\sigma^2(1-\rho)-\sigma^2-\mu^2 = \sigma^2(1-2\rho)-\mu^2$, which is negative if and only if $\rho > \tfrac12 - \mu^2/(2\sigma^2)$.

\emph{(iii)} With $m(\hat x_{t+1})=\mu$, $d_{t+1}=m_{t+1}-\mu$, $\E[d^2]=\sigma^2$, and the misspecification term is $-2\mu\E[u]=0$. Hence $\MSE(\hat y)-\MSE(\pers) = \sigma^2-(\sigma^2+\mu^2) = -\mu^2$. Subtracting the expression in (ii) gives $\MSE(\text{climatological})-\MSE(\text{persisted}) = -\mu^2 - \sigma^2(1-2\rho)+\mu^2 = -\sigma^2(1-2\rho)$. \qed

\paragraph{Derivation of Remark~\ref{rem:crossover}.} For weakly stationary $y$ with mean $\bar y$, variance $\sigma_y^2$ and lag-$h$ autocorrelation $\rho_h$, the $h$-step persistence error is $y_{t+h}-y_t$ with mean zero and variance $2\sigma_y^2(1-\rho_h)$, and the error of the known-mean forecast is $y_{t+h}-\bar y$ with mean square $\sigma_y^2$. The ratio of root mean squared errors is therefore $\sqrt{2(1-\rho_h)}$, which is below one if and only if $\rho_h>\tfrac12$. With the mean estimated from $n_t$ training values, $\MSE_h(\mathrm{mean}) = \sigma_y^2 + \operatorname{Var}(\hat{\bar y}_t) - 2\operatorname{cov}(y_{t+h},\hat{\bar y}_t)$, whose correction terms are of order $\sigma_y^2/n_t$ times sums of autocorrelations and vanish as $n_t\to\infty$. \qed

\section*{Declarations}

\paragraph{Lineage} This article unifies two single-system studies, a scored surplus-production ladder on Northern cod \citep{Abaee2026cod} and a scored one-pool water-balance ladder on Edwards Aquifer head \citep{Abaee2026edwards}, whose archived scores it reports unchanged, and it supersedes their earlier two-part combination. The two systems are reported under their own protocols and are not pooled.

\paragraph{Data availability} All input data, analysis scripts, result files and the frozen protocol documents of both systems are archived at \url{https://github.com/MIKEAA2020/general-sustainability}, with Zenodo records for the cod study (\url{https://zenodo.org/records/22553609}) and the aquifer study (\url{https://doi.org/10.5281/zenodo.22552680}). Cod: the NCAM Table A2 and xteNCAM Table 17 files, the landings tables, the capelin acoustic index and the RAM Legacy Stock Assessment Database v4.66 extract (10.5281/zenodo.14043038; stock COD2J3KL) are archived in-repository under \texttt{paperE1\_calibration\_data\_v1\_wave\_e\_cod/} and \texttt{wave\_e\_cod/data/} with their source record. Aquifer: J-17 daily highs are Texas Water Development Board well 6837203; recharge is the USGS data release of \citet{UmphresChoi2025}; pumpage is Edwards Aquifer Authority Table 1 \citep{EAA2024}; Comal Springs discharge is USGS 08168710; Ni\~no~3.4 is NOAA PSL HadISST (raw file registered); precipitation is NCEI nClimDiv, whose raw file is not redistributed (provenance URL archived), so the three precipitation columns of the registered twenty-column panel \texttt{wave\_e\_edwards/data/annual\_panel.csv} rebuild from the provenance source rather than from the archive alone, while scoring from the registered panel does not require it.

\paragraph{Code availability} The cod ladder and its campaigns are produced by \texttt{wave\_e\_cod/src/run\_ladder.py}, \texttt{run\_xte.py}, \texttt{run\_capelin\_regime.py}, \texttt{run\_capelin\_index.py}, the uncertainty script \texttt{campaign\_e1\_dm\_uncertainty\_regime.py} and the simulation \texttt{sim\_retention\_power.py}, with the paired bootstrap, period decomposition, catch-timing sensitivity and figure scripts archived under the previous edition's \texttt{work/} directory and carried in \texttt{paper11\_forecasting\_baselines\_v68\_src/work/}. The aquifer ladder is produced by \texttt{wave\_e\_edwards/src/run\_ladder.py}, the climate pass by \texttt{run\_recharge.py}, the uncertainty layer by \texttt{e3\_audit\_uncertainty.py} with its independent replication \texttt{campaign\_e3\_dm\_uncertainty.py}, and the pumpage counterfactuals by \texttt{campaign\_e3\_pumpage\_scenarios.py}; the verification script \texttt{paperE3\_edwards\_forecast\_ladder\_v17\_verification.py} re-runs the ladder from the registered panel and recompares every results-table cell. The new cross-system analyses are the baseline horizon profile and full-sample autocorrelation checks, produced by \texttt{work/horizon\_profile.py}; Fig.~\ref{fig:mechanism} is produced by \texttt{work/make\_fig\_mechanism.py}. The exact one-step J-17 decomposition is reproduced by \texttt{work/edw/increment\_identity\_rolling.py}, which uses each origin's fitted full M2 increment, checks the identity at every origin, and matches clipped M2 and persistence forecasts to the archived file; its per-origin CSV and summary JSON are archived with the source bundle, with the summary reported in SI S23. The deterministic cod operating-characteristic rerun is produced by \texttt{work/cod/sim\_retention\_power\_repro.py}, with the batch aggregator \texttt{aggregate\_power\_simulation\_repro.py}. The raw output, summary tables and combined run manifest are archived with the script only after the complete replicate-ID range has passed coverage and hash checks; partial batch progress is not reported as a result. The standard-library verifier \texttt{paper11\_forecasting\_baselines\_v68\_verification.py} checks the printed values, the 75-origin identity and archived forecast agreement, and textual invariants.

\paragraph{Reproducibility} All computations are deterministic within a fixed interpreter and library stack. The cod runners were re-executed in a third software environment for the previous edition and every scored row reproduced at printed precision, with one known environment sensitivity (the M1b optimum on Specification B, $\pm17\kt$, immaterial to the verdict; SI S13). The aquifer scripts regenerate every archived result file byte for byte in a fresh environment, and the two implementations of the uncertainty layer agree on every conclusion under different seeds, block lengths and HAC conventions (SI S18).

\paragraph{CRediT authorship contribution statement} Amin Abaee: Conceptualization, Methodology, Software, Validation, Formal analysis, Investigation, Data curation, Writing -- original draft, Writing -- review \& editing, Visualization.

\paragraph{Funding} This research did not receive any specific grant from funding agencies in the public, commercial, or not-for-profit sectors.

\paragraph{Declaration of competing interest} The author declares no known competing financial interests or personal relationships that could have appeared to influence the work reported in this paper.

\paragraph{Declaration of generative AI in the writing process} GLM (Z.ai), Qwen (Alibaba Cloud) and DeepSeek AI assisted with drafting and iterative review. The author reviewed and edited all outputs and takes full responsibility for the final work.

\begin{thebibliography}{99}
\providecommand{\natexlab}[1]{#1}

\bibitem[Abaee(2026a)]{Abaee2026edwards}
Abaee, A., 2026a. Does a one-pool water-balance model improve forecasts of Edwards Aquifer head? A scored test at J-17. Zenodo. \url{https://doi.org/10.5281/zenodo.22552680}

\bibitem[Abaee(2026b)]{Abaee2026cod}
Abaee, A., 2026b. Does a surplus-production ladder improve forecasts of Northern cod? A scored retention test against persistence on NAFO 2J3KL. Zenodo. \url{https://zenodo.org/records/22553609}

\bibitem[Adamowski and Chan(2011)]{AdamowskiChan2011}
Adamowski, J., Chan, H.F., 2011. A wavelet neural network conjunction model for groundwater level forecasting. J. Hydrol. 407, 28--40. \url{https://doi.org/10.1016/j.jhydrol.2011.06.013}

\bibitem[Bergmeir and Ben{\'\i}tez(2012)]{Bergmeir2012}
Bergmeir, C., Ben{\'\i}tez, J.M., 2012. On the use of cross-validation for time series predictor evaluation. Inf. Sci. 191, 192--213.

\bibitem[Buren et al.(2014)]{Buren2014}
Buren, A.D., Koen-Alonso, M., Pepin, P., Mowbray, F., Nakashima, B., Stenson, G., Ollerhead, N., Montevecchi, W.A., 2014. Bottom-up regulation of capelin, a keystone forage species. PLoS ONE 9, e87589.

\bibitem[Cadigan(2016)]{Cadigan2016}
Cadigan, N.G., 2016. A state-space stock assessment model for northern cod, including under-reported catches and variable natural mortality rates. Can. J. Fish. Aquat. Sci. 73, 296--308.

\bibitem[Carvalho et al.(2021)]{Carvalho2021}
Carvalho, F., Winker, H., Courtney, D., Kapur, M., Kell, L., Cardinale, M., Schirripa, M., Kitakado, T., Yemane, D., Piner, K.R., Maunder, M.N., Taylor, I., Wetzel, C.R., Doering, K., Johnson, K.F., Methot, R.D., 2021. A cookbook for using model diagnostics in integrated stock assessments. Fish. Res. 240, 105959.

\bibitem[Daliakopoulos et al.(2005)]{Daliakopoulos2005}
Daliakopoulos, I.N., Coulibaly, P., Tsanis, I.K., 2005. Groundwater level forecasting using artificial neural networks. J. Hydrol. 309, 229--240. \url{https://doi.org/10.1016/j.jhydrol.2004.12.001}

\bibitem[DFO(2009)]{DFO2009}
DFO, 2009. A fishery decision-making framework incorporating the Precautionary Approach. Fisheries and Oceans Canada, Ottawa.

\bibitem[DFO(2016)]{DFO2016}
DFO, 2016. Stock assessment of Northern cod (NAFO Divs. 2J3KL) in 2016. DFO Can. Sci. Advis. Sec. Sci. Advis. Rep. 2016/026.

\bibitem[DFO(2024a)]{DFO2024a}
DFO, 2024a. NAFO Divisions 2J3KL Northern Cod stock assessment to 2024. DFO Can. Sci. Advis. Sec. Sci. Advis. Rep. 2024/049.

\bibitem[DFO(2024b)]{DFO2024b}
DFO, 2024b. Assessment of capelin in NAFO Divisions 2J3KL. DFO Can. Sci. Advis. Sec. Sci. Advis. Rep. 2024/050.

\bibitem[Diebold and Mariano(1995)]{Diebold1995}
Diebold, F.X., Mariano, R.S., 1995. Comparing predictive accuracy. J. Bus. Econ. Stat. 13, 253--263.

\bibitem[Diebold(2015)]{Diebold2015}
Diebold, F.X., 2015. Comparing predictive accuracy, twenty years later: a personal perspective on the use and abuse of Diebold--Mariano tests. J. Bus. Econ. Stat. 33, 1--9.

\bibitem[Edwards Aquifer Authority(2024)]{EAA2024}
Edwards Aquifer Authority, 2024. 2023 Groundwater Discharge and Usage, Table 1 (after USGS letter report, 5 April 2024). Edwards Aquifer Authority, San Antonio, TX. \url{https://www.edwardsaquifer.org/wp-content/uploads/2025/05/2023-Groundwater-Discharge-and-Usage.pdf}

\bibitem[Harvey et al.(1997)]{Harvey1997}
Harvey, D., Leybourne, S., Newbold, P., 1997. Testing the equality of prediction mean squared errors. Int. J. Forecast. 13, 281--291.

\bibitem[Hilborn and Walters(1992)]{HilbornWalters1992}
Hilborn, R., Walters, C.J., 1992. Quantitative Fisheries Stock Assessment: Choice, Dynamics and Uncertainty. Chapman and Hall, New York.

\bibitem[Hutchings and Myers(1994)]{Hutchings1994}
Hutchings, J.A., Myers, R.A., 1994. What can be learned from the collapse of a renewable resource? Atlantic cod, \emph{Gadus morhua}, of Newfoundland and Labrador. Can. J. Fish. Aquat. Sci. 51, 2126--2146.

\bibitem[Hutchings(1996)]{Hutchings1996}
Hutchings, J.A., 1996. Spatial and temporal variation in the density of northern cod and a review of hypotheses for the stock's collapse. Can. J. Fish. Aquat. Sci. 53, 943--962.

\bibitem[Hutchings(2005)]{Hutchings2005}
Hutchings, J.A., 2005. Life history consequences of overexploitation to population recovery in Northwest Atlantic cod (\emph{Gadus morhua}). Can. J. Fish. Aquat. Sci. 62, 824--832.

\bibitem[Hyndman and Athanasopoulos(2021)]{HyndmanAthanasopoulos2021}
Hyndman, R.J., Athanasopoulos, G., 2021. Forecasting: Principles and Practice, third ed. OTexts, Melbourne.

\bibitem[Hyndman and Koehler(2006)]{HyndmanKoehler2006}
Hyndman, R.J., Koehler, A.B., 2006. Another look at measures of forecast accuracy. Int. J. Forecast. 22, 679--688.

\bibitem[Kell et al.(2016)]{Kell2016}
Kell, L.T., Kimoto, A., Kitakado, T., 2016. Evaluation of the prediction skill of stock assessment using hindcasting. Fish. Res. 183, 119--127.

\bibitem[Kell et al.(2021)]{Kell2021}
Kell, L.T., Sharma, R., Kitakado, T., Winker, H., Mosqueira, I., Cardinale, M., Fu, D., 2021. Validation of stock assessment methods: is it me or my model talking? ICES J. Mar. Sci. 78, 2244--2255.

\bibitem[Kokkalis et al.(2024)]{Kokkalis2024}
Kokkalis, A., Berg, C.W., Kapur, M.S., Winker, H., Jacobsen, N.S., Taylor, M.H., Ichinokawa, M., Miyagawa, M., Medeiros-Leal, W., Nielsen, J.R., Mildenberger, T.K., 2024. Good practices for surplus production models. Fish. Res. 275, 107010.

\bibitem[K{\"u}nsch(1989)]{Kunsch1989}
K{\"u}nsch, H.R., 1989. The jackknife and the bootstrap for general stationary observations. Ann. Stat. 17, 1217--1241.

\bibitem[Liermann and Hilborn(1997)]{Liermann1997}
Liermann, M., Hilborn, R., 1997. Depensation in fish stocks: a hierarchic Bayesian meta-analysis. Can. J. Fish. Aquat. Sci. 54, 1976--1984.

\bibitem[Ludwig and Walters(1985)]{LudwigWalters1985}
Ludwig, D., Walters, C.J., 1985. Are age-structured models appropriate for catch-effort data? Can. J. Fish. Aquat. Sci. 42, 1066--1072.

\bibitem[Makridakis and Hibon(2000)]{Makridakis2000}
Makridakis, S., Hibon, M., 2000. The M3-Competition: results, conclusions and implications. Int. J. Forecast. 16, 451--476.

\bibitem[Makridakis et al.(2018)]{Makridakis2018}
Makridakis, S., Spiliotis, E., Assimakopoulos, V., 2018. The M4 Competition: results, findings, conclusion and way forward. Int. J. Forecast. 34, 802--808.

\bibitem[Maunder and Piner(2015)]{Maunder2015}
Maunder, M.N., Piner, K.R., 2015. Contemporary fisheries stock assessment: many issues still remain. ICES J. Mar. Sci. 72, 7--18.

\bibitem[Murphy et al.(2025)]{Murphy2025}
Murphy, H.M., Adamack, A.T., Lewis, R.S., Bourne, C.M., 2025. Assessment of capelin in NAFO Divisions 2J+3KL to 2023. DFO Can. Sci. Advis. Sec. Res. Doc. 2025/022.

\bibitem[Myers et al.(1995)]{Myers1995}
Myers, R.A., Barrowman, N.J., Hutchings, J.A., Rosenberg, A.A., 1995. Population dynamics of exploited fish stocks at low population levels. Science 269, 1106--1108.

\bibitem[Myers et al.(1997)]{Myers1997}
Myers, R.A., Hutchings, J.A., Barrowman, N.J., 1997. Why do fish stocks collapse? The example of cod in Atlantic Canada. Ecol. Appl. 7, 91--106.

\bibitem[Northwest Atlantic Fisheries Centre(2025)]{NAFC2025}
Northwest Atlantic Fisheries Centre, 2025. 2J3KL cod and capelin biomass indices. Zenodo. \url{https://doi.org/10.5281/zenodo.17515115}

\bibitem[Ohmer et al.(2026)]{Ohmer2026}
Ohmer, M., Liesch, T., Habbel, B., Heudorfer, B., Gomez, M., Clos, P., N{\"o}lscher, M., Broda, S., 2026. GEMS-GER: a machine learning benchmark dataset of long-term groundwater levels in Germany with meteorological forcings and site-specific environmental features. Earth Syst. Sci. Data 18, 77.

\bibitem[Pedersen and Berg(2017)]{PedersenBerg2017}
Pedersen, M.W., Berg, C.W., 2017. A stochastic surplus production model in continuous time. Fish Fish. 18, 226--243.

\bibitem[Pella and Tomlinson(1969)]{Pella1969}
Pella, J.J., Tomlinson, P.K., 1969. A generalized stock production model. Bull. Inter-Am. Trop. Tuna Comm. 13, 419--496.

\bibitem[Perretti et al.(2013)]{Perretti2013}
Perretti, C.T., Munch, S.B., Sugihara, G., 2013. Model-free forecasting outperforms the correct mechanistic model for simulated and experimental data. Proc. Natl. Acad. Sci. U.S.A. 110, 5253--5257.

\bibitem[Polacheck et al.(1993)]{Polacheck1993}
Polacheck, T., Hilborn, R., Punt, A.E., 1993. Fitting surplus production models: comparing methods and measuring uncertainty. Can. J. Fish. Aquat. Sci. 50, 2597--2607.

\bibitem[Puente(1978)]{Puente1978}
Puente, C., 1978. Method of estimating natural recharge to the Edwards Aquifer in the San Antonio area, Texas. U.S. Geological Survey Water-Resources Investigations Report 78-10, Austin, TX.

\bibitem[Punt et al.(2016)]{Punt2016}
Punt, A.E., Butterworth, D.S., de Moor, C.L., De Oliveira, J.A.A., Haddon, M., 2016. Management strategy evaluation: best practices. Fish Fish. 17, 303--334.

\bibitem[Punt(2003)]{Punt2003}
Punt, A.E., 2003. Extending production models to include process error in the population dynamics. Can. J. Fish. Aquat. Sci. 60, 1217--1228.

\bibitem[Regular et al.(2025)]{Regular2025}
Regular, P.M., et al., 2025. Assessment of the Northern Cod stock in NAFO Divisions 2J3KL in 2024. DFO Can. Sci. Advis. Sec. Res. Doc. 2025/048.

\bibitem[Ricard et al.(2012)]{Ricard2012}
Ricard, D., Minto, C., Jensen, O.P., Baum, J.K., 2012. Examining the knowledge base and status of commercially exploited marine species with the RAM Legacy Stock Assessment Database. Fish Fish. 13, 380--398.

\bibitem[Romano and Wolf(2005)]{RomanoWolf2005}
Romano, J.P., Wolf, M., 2005. Stepwise multiple testing as formalized data snooping. Econometrica 73, 1237--1282.

\bibitem[Ropelewski and Halpert(1986)]{Ropelewski1986}
Ropelewski, C.F., Halpert, M.S., 1986. North American precipitation and temperature patterns associated with the El Ni{\~n}o/Southern Oscillation (ENSO). Mon. Weather Rev. 114, 2352--2362.

\bibitem[Rose and Rowe(2015)]{RoseRowe2015}
Rose, G.A., Rowe, S., 2015. Northern cod comeback. Can. J. Fish. Aquat. Sci. 72, 1789--1798.

\bibitem[Rose and Walters(2019)]{RoseWalters2019}
Rose, G.A., Walters, C.J., 2019. The state of Canada's iconic Northern cod: a second opinion. Fish. Res. 219, 105314.

\bibitem[Rose(2026)]{Rose2026}
Rose, G.A., 2026. Northern cod comeback: 10 years after. Can. J. Fish. Aquat. Sci. 83, 1--14. \url{https://doi.org/10.1139/cjfas-2025-0141}

\bibitem[Scanlon et al.(2003)]{Scanlon2003}
Scanlon, B.R., Mace, R.E., Barrett, M.E., Smith, B., 2003. Can we simulate regional groundwater flow in a karst system using equivalent porous media models? Case study, Barton Springs Edwards aquifer, USA. J. Hydrol. 276, 137--158.

\bibitem[Schaefer(1954)]{Schaefer1954}
Schaefer, M.B., 1954. Some aspects of the dynamics of populations important to the management of the commercial marine fisheries. Bull. Inter-Am. Trop. Tuna Comm. 1(2), 27--56.

\bibitem[Schijns et al.(2021)]{Schijns2021}
Schijns, R., Froese, R., Hutchings, J.A., Pauly, D., 2021. Five centuries of cod catches in Eastern Canada. ICES J. Mar. Sci. 78, 2675--2683.

\bibitem[Shelton and Healey(1999)]{Shelton1999}
Shelton, P.A., Healey, B.P., 1999. Should depensation be dismissed as a possible explanation for the lack of recovery of the northern cod (\emph{Gadus morhua}) stock? Can. J. Fish. Aquat. Sci. 56, 1521--1524.

\bibitem[Tashman(2000)]{Tashman2000}
Tashman, L.J., 2000. Out-of-sample tests of forecasting accuracy: an analysis and review. Int. J. Forecast. 16, 437--450.

\bibitem[Umphres and Choi(2025)]{UmphresChoi2025}
Umphres, G.D., Choi, N.J., 2025. Estimated annual recharge to the Edwards Aquifer in the San Antonio area, by stream basin or ungaged area, 1934--2024. U.S. Geological Survey data release. \url{https://doi.org/10.5066/P1BI62NY}

\bibitem[Walters and Maguire(1996)]{Walters1996}
Walters, C., Maguire, J.-J., 1996. Lessons for stock assessment from the northern cod collapse. Rev. Fish Biol. Fish. 6, 125--137.

\bibitem[Ward et al.(2014)]{Ward2014}
Ward, E.J., Holmes, E.E., Thorson, J.T., Collen, B., 2014. Complexity is costly: a meta-analysis of parametric and non-parametric methods for short-term population forecasting. Oikos 123, 652--661.

\bibitem[White(2000)]{White2000}
White, H., 2000. A reality check for data snooping. Econometrica 68, 1097--1126.

\bibitem[Zhu et al.(2026)]{Zhu2026}
Zhu, Q., Zhu, Y., Niu, J., Huang, J., Huang, F., Zhou, X., Liu, D., Hu, B.X., 2026. Benchmarking machine learning and deep learning models for groundwater level prediction in karst aquifers: the dominant role of hydrogeological complexity. Water 18, 939. \url{https://doi.org/10.3390/w18080939}

\end{thebibliography}

\end{document}
